% Fragmento público generado desde propietarios íntegros.
% Residencia: 52.9.2.
% Fuente: ../incoming/radion_monografia_20260827/manuscrito/sections/10_minkowski_hodge_lorentz.tex:101-186
\subsubsection{Realización lorentziana de las trayectorias helicoidales}

En un campo magnético uniforme, una partícula de carga \(q\), masa \(m\) y
factor relativista \(\gamma\) posee
\begin{equation}
\omega_c=\frac{|q|B}{\gamma m},
\qquad
r_L=\frac{p_\perp}{|q|B},
\qquad
\Delta z=\frac{2\pi p_\parallel}{|q|B}.
\label{eq:lorentz-helix}
\end{equation}
La fase ciclotrónica cierra mientras el movimiento paralelo avanza: la órbita
es helicoidal. El paralelismo con \eqref{eq:H9-helix} es estructural y puede
formalizarse por un transductor que preserve fase, orientación y paso. No se
identifica el cociente \(q_9\) con la coordenada espacial \(z\); se reconoce
la misma arquitectura de retorno más avance.

La fuerza de Lorentz no se deriva de la imagen de una espiral por semejanza.
Se deriva al completar \eqref{eq:lorentz-chain}. La imagen ayuda a comprender
por qué fase circular y memoria longitudinal son compatibles; la forma de
campo determina la dinámica.

\subsubsection{Confluencia entre acción, área, información y gravedad}

\paragraph{Realización Cartan--Holst de la holonomía}

La promoción gravitatoria se expresa mediante un morfismo de holonomías
\begin{equation}
\operatorname{Hol}_{\mathcal A}
\bigl(\mathfrak R_{\mathrm{grav}}(\gamma)\bigr)
=\rho_C\bigl(\operatorname{Hol}_{\TPK}(\gamma)\bigr).
\label{eq:holonomy-gravity}
\end{equation}
El lado derecho contiene la memoria de ruta; \(\rho_C\) la realiza en una
conexión gravitatoria. La torsión de Cartan
\[
T^I=D_\omega e^I
\]
aparece en el codominio. La torsión global \(\Delta_4\) es una coordenada
upstream que puede alimentar el mapa, pero no se renombra como \(T^I\) sin
\eqref{eq:holonomy-gravity}.

\paragraph{Transducción del área de acción en área informacional}

El radión fija un cuanto de acción por fase. El lector hexada--octada de
Barbero usa posteriormente los canales \(q_\pm\), la incidencia \(90/120\) y
el funcional \eqref{eq:barbero-núcleo}. La dirección causal es
\[
(A,C^*,\hret)\to q_\pm\to K_{6\to8}
\to\gamma_{\mathrm{BI}}\to\text{área/información}.
\]
El alfabeto trítico aporta la unidad de entropía \(\log3\) y un corte de
\(81\) canales aporta \(81\log3\). Esta entropía no se confunde con
\(\log2\), con \(\log\varphi\), con una entropía microcanónica ni con una
entropía de von Neumann sin especificar el estado.

La síntesis física es afirmativa pero tipada: el radión aporta la unidad de
acción orientada que una holonomía puede transportar; Barbero publica la
normalización de área; la pantalla causal y Hodge permiten realizar campos;
Cartan--Holst recibe la holonomía en geometría gravitatoria.

\begin{statusbox}
La ecuación escalar del radión no es por sí sola una ecuación de Einstein. Su
aportación a la gravedad es ser una sección exacta de la arquitectura de
acción, orientación y holonomía que la realización gravitatoria consume.
\end{statusbox}

\subsubsection{Transformación de Wick como cambio de lector analítico}

En la formulación convencional, la rotación de Wick relaciona una variable
temporal lorentziana con una coordenada euclídea mediante una continuación
compleja. En HMT, la raíz \(J^2=-I\), el involutivo \(R^2=I\) y el umbral
\(N^2=0\) ya existen antes. El cambio de lector
\[
J\longleftrightarrow R
\]
transporta una pareja común/orientada entre publicación elíptica y
publicación hiperbólica. La continuación compleja reconoce después esa
operación en coordenadas analíticas.

Esto explica por qué círculo, elipse de acción e interior hiperbólico pueden
pertenecer al mismo estado sin ser la misma métrica. El círculo usa \(J\) para
la fase; la forma de la elipse usa un compresión simpléctica generado por \(R\); Klein mide
el interior mediante razón doble; la pantalla de Minkowski realiza causalidad
en otro codominio.
